\documentclass[11pt,a4paper]{article}
\usepackage[a4paper,margin=2.6cm]{geometry}
\usepackage{iftex}
\ifPDFTeX
  \usepackage[T2A]{fontenc}
  \usepackage[utf8]{inputenc}
  \usepackage[russian]{babel}
\else
  \usepackage{fontspec}
  \usepackage{polyglossia}
  \setdefaultlanguage{russian}
  \setotherlanguage{english}
  \defaultfontfeatures{Ligatures=TeX}
  % Цепочка шрифтов: Times New Roman -> Liberation (Overleaf) -> CMU.
  \IfFontExistsTF{Times New Roman}{%
    \setmainfont{Times New Roman}%
    \setsansfont{Arial}%
    \setmonofont{Courier New}%
    \newfontfamily\cyrillicfont{Times New Roman}%
    \newfontfamily\cyrillicfontsf{Arial}%
    \newfontfamily\cyrillicfonttt{Courier New}%
  }{%
    \IfFontExistsTF{Liberation Serif}{%
      \setmainfont{Liberation Serif}%
      \setsansfont{Liberation Sans}%
      \setmonofont{Liberation Mono}%
      \newfontfamily\cyrillicfont{Liberation Serif}%
      \newfontfamily\cyrillicfontsf{Liberation Sans}%
      \newfontfamily\cyrillicfonttt{Liberation Mono}%
    }{%
      \setmainfont{CMU Serif}%
      \setsansfont{CMU Sans Serif}%
      \setmonofont{CMU Typewriter Text}%
      \newfontfamily\cyrillicfont{CMU Serif}%
      \newfontfamily\cyrillicfontsf{CMU Sans Serif}%
      \newfontfamily\cyrillicfonttt{CMU Typewriter Text}%
    }%
  }
\fi
\usepackage{microtype}
\usepackage{mathtools,amssymb}
\usepackage{graphicx}
\usepackage{booktabs}
\usepackage{array}
\usepackage{float}
\usepackage{caption}
\captionsetup{font=small,labelfont=bf}
\usepackage{enumitem}
\setlist{nosep,leftmargin=1.5em}
\usepackage{csquotes}
\usepackage[backend=biber,style=gost-numeric,sorting=none]{biblatex}
\addbibresource{references_article1.bib}
\usepackage[unicode,hidelinks]{hyperref}
\setlength{\emergencystretch}{2em}

\newcommand{\Psurr}{P_{\mathrm{surr}}}

\title{{\normalsize\normalfont\hspace*{-\leftskip}\parbox{\linewidth}{\raggedright УДК~551.513}}\\[1.5em]
Сопряжённость иерархических уровней атмосферной циркуляции\\
как региональный инвариант}
\author{Н.~Сотириади\thanks{[ЗАПОЛНИТЬ ПЕРЕД ПОДАЧЕЙ] организация, почтовый адрес с индексом, страна. E-mail: \texttt{n.sotiriadi@gmail.com}. ORCID: [ЗАПОЛНИТЬ].}%
\thanks{Первая статья серии из трёх. Протоколы, код и численные результаты всех работ серии открыты: \url{https://github.com/Theclimateguy/Shroedinger} (архив версий --- Zenodo, DOI 10.5281/zenodo.19565769).}}
\date{2026}

\begin{document}
\maketitle

\begin{abstract}
Спектральный и фрактальный анализ географических полей отвечает на вопрос, на каких масштабах располагаются уровни организации геосистемы. Следующий вопрос касается отношения между уровнями: предопределено ли размещение мезомасштабной изменчивости соседним крупным уровнем или она размещается по собственным локальным причинам. В работе по полю относительного вихря на поверхности 850~гПа (реанализ ERA5, 2017--2024~гг.; 12 основных и 27 контрольных областей) построена мера такого отношения --- профиль сопряжённости соседних октавных уровней $P$: пространственная корреляция карт активности двух соседних масштабов, очищенная от вклада самой процедуры разложения. Проверена гипотеза о том, что профиль является региональным инвариантом в операциональном смысле: сохраняет своё значение при смене сезона, года, фазы ЭНСО, источника данных и схемы разбиения территории, различая при этом циркуляционные режимы. Гипотеза подтверждена на выборках, не участвовавших в построении величины ($p=0{,}001$). Четыре альтернативных объяснения --- численная диссипация реанализа, спектр мощности поля, геометрия градусной сетки и устройство наблюдательной сети --- проверены прямыми контролями: региональная сигнатура сохраняется при ограничении разрешёнными ступенями, содержит воспроизводимую информацию сверх спектра (контроль равновеликими по площади областями: $p=0{,}033$), воспроизводится при равновеликой нарезке территории и в свободном модельном расчёте, не усваивающем атмосферных наблюдений (ранговая корреляция 0,71 при внутреннем пределе метода 0,93). Тенденция географии --- более высокие значения над внетропическими штормовыми путями и более низкие над влажно-конвективными режимами --- отчётлива на тонких полосах и сохраняется с исключениями на разрешённых ступенях; бассейны Конго и Амазонии при сходных стандартных характеристиках режима различаются архитектурой иерархии. Величина дополняет аппарат районирования характеристикой устройства межуровневых связей; границы её применимости, включая отвергнутые интерпретации, перечислены явно.

\smallskip
\noindent\textbf{Ключевые слова:} иерархическая организация геосистем, инварианты динамической геосистемы, масштаб, атмосферная циркуляция, штормовые пути, реанализ, районирование.
\end{abstract}

\section{Постановка вопроса}
\label{sec:intro}

\subsection{Иерархическая организация и инварианты территории}

Представление о геосистеме как об открытой, динамической и иерархически организованной целостности, в которой различаются топологический, региональный и планетарный порядки размерности, составляет основание учения о геосистемах В.\,Б.~Сочавы \cite{Sochava1978}. Независимо близкая рамка сложилась в западноевропейской физической географии: Ж.~Бертран связал ландшафтные единицы с пространственно-временным масштабом и предложил их вложенную классификацию --- от природной зоны до геотопа \cite{Bertrand1968}. Общесистемное основание такой архитектуры дано Г.~Саймоном: сложные системы, как правило, иерархичны и <<почти разложимы>> --- быстрые процессы внутри подсистем относительно автономны, тогда как более медленные межуровневые связи действуют агрегированно \cite{Simon1962}. Иерархичность геосистем, таким образом, представляет собой общую архитектуру сложной организации, разделяемую несколькими научными школами.

Из этой архитектуры вытекает принципиальная методологическая трудность, наиболее ясно сформулированная в экологии: у наблюдаемого пространственного рисунка нет единственного <<естественного>> масштаба описания, а механизмы и порождаемые ими паттерны могут относиться к разным масштабам, поэтому предметом исследования должна становиться сама перестройка организации поля при смене масштаба \cite{Levin1992}; применительно к пространственно неоднородной среде эта постановка конкретизирована в концепции иерархической динамики мозаик \cite{WuLoucks1995}. Отсюда эмпирическая задача, определяющая настоящую работу: строить воспроизводимые характеристики того, как организация изменчивости соотносится между явно заданными диапазонами масштаба, и проверять, в каких географических условиях эти характеристики сохраняются.

В отечественной географии эта постановка получила операциональное развитие. В работах Ю.\,Г.~Пузаченко пространственно-временная иерархия геосистем рассмотрена с позиций теории колебаний \cite{Puzachenko1986}, сформирован аппарат выделения уровней организации по спектральным и фрактальным свойствам географических полей \cite{Puzachenko2004}, а поиск \emph{инвариантов динамики геосистем} --- устойчивых во времени числовых характеристик организации, различающих территории, --- поставлен как самостоятельная исследовательская программа \cite{Puzachenko2010}. Инвариант в этом понимании представляет собой воспроизводимую меру устройства территории, сохраняющую значение при смене состояния системы и потому пригодную для сравнения и районирования. Непосредственным развитием линии стала работа А.\,Н.~Кренке с соавторами \cite{Krenke2019}, где иерархические уровни организации регионального мезоклимата выделялись по двумерным спектрам главных компонент климатических полей. Тем самым получен ответ на вопрос, \emph{где} расположены уровни организации.

\subsection{Нерешённый вопрос}

Выделение уровней --- операция классификационная; следующий по логике вопрос касается \emph{отношения} между уровнями. Мезомасштабная изменчивость --- вихри и конвективные структуры размером в сотни километров --- может размещаться там, где активен соседний крупный уровень, и тогда крупное поле предопределяет её географию. Она может размещаться и независимо от крупного поля, по собственным локальным причинам. Между этими крайностями лежит континуум, и положение конкретной территории на нём до сих пор не измерялось: величины такого рода не строились как самостоятельные характеристики территории и не проверялись на статус инвариантов. Настоящая работа посвящена этому вопросу применительно к атмосферной циркуляции.

Выбор объекта не случаен и в физическом отношении. Диапазон масштабов работы (50--1600~км) охватывает известный переход в спектре кинетической энергии атмосферы --- от крупномасштабной ветви с наклоном, близким к $k^{-3}$, к мезомасштабной ветви, близкой к $k^{-5/3}$ \cite{NastromGage1985}: сопряжённость соседних масштабных полос измеряется здесь поперёк того диапазона, где сама форма спектра свидетельствует о смене режима организации.

Близкие по духу инструменты развивались в трёх традициях, и каждая оставляет искомую величину за пределами рассмотрения. Каскадные и мультифрактальные описания \cite{LovejoySchertzer2013,Lovejoy2011,Kouhen2024} характеризуют распределение интенсивности \emph{вдоль} масштабного континуума; пространственная согласованность карт активности \emph{между} конкретными уровнями остаётся вне их предмета. Вейвлет-огибающие и их статистики --- коэффициент амплитудной модуляции \cite{Farge1992,Mathis2009} и систематизирующие его скаттеринг-ковариации \cite{Mallat2012,BrunaMallat2013,Cheng2024} --- измеряют факт связи уровней, но применялись как признаки для классификации; одно из недавних приложений подобной техники к различению режимов геофизических полей --- океанологическая работа \cite{Skinner2025}. Литература о предсказуемости атмосферы \cite{Lorenz1969,Zhang2007,RotunnoSnyder2008,Judt2020,SelzCraig2023} отводит межмасштабному взаимодействию роль канала, по которому неопределённость мелких масштабов заражает синоптический уровень, и моделирует силу этой связи; измерение этой связи по данным остаётся за пределами названных работ. Сама форма оценки --- корреляция огибающих соседних полос одного поля --- известна: её ближайший формальный родственник введён в нейрофизиологии для выявления связи некогерентных компонент сигнала \cite{Bruns2000}; в той же дисциплине изучается более широкий класс межчастотных связей, включая фазо-амплитудную \cite{Canolty2006}; родственная форма --- корреляция крупномасштабного сигнала с огибающей мелких масштабов --- независимо сложилась в физике пристеночной турбулентности \cite{Mathis2009}. Независимое появление родственных форм в разных дисциплинах подчёркивает естественность конструкции и оставляет открытыми два вопроса: о её содержании для конкретного географического поля и о статусе инварианта.

\subsection{Гипотеза}
\label{sec:hyp}

Гипотеза настоящей работы состоит в том, что сопряжённость соседних уровней циркуляции является устойчивым региональным свойством: она различает циркуляционные режимы и сохраняет своё значение при смене сезона, года и фазы климатической изменчивости. В терминологии программы поиска инвариантов динамики геосистем \cite{Puzachenko2010} такое свойство называется региональным инвариантом. Слово требует точного определения. Пусть $P^{(\ell)}_{G;\tau}$ --- значение профиля сопряжённости области $G$ на масштабной ступени $\ell$ при условиях наблюдения $\tau=(s,y,e,d,g)$: сезон $s$, год $y$, фаза ЭНСО (Эль-Ниньо --- Южное колебание) $e$, источник данных $d$, схема разбиения территории $g$. Проверяемое утверждение --- условная инвариантность на заранее заявленном классе преобразований $\mathcal{T}$:
\begin{equation}
\label{eq:inv}
P^{(\ell)}_{G;\tau}\approx P^{(\ell)}_{G;\tau'}\quad\text{для всех }\tau,\tau'\in\mathcal{T},
\end{equation}
при том что различия между областями остаются существенно больше различий внутри области (форма утверждения в терминах кластеризации --- разд.~\ref{sec:test}). Класс $\mathcal{T}$ зафиксирован протоколами до вычислений и включает смену сезона (январь--март / июль--сентябрь), смену независимого года, смену фазы ЭНСО, смену источника данных (ERA5 / MERRA-2 / свободный модельный расчёт) и смену схемы разбиения территории (градусные / равновеликие области).

Столь же явно фиксируются границы утверждения. Оно относится к полю относительного вихря; к уровню 850~гПа (перенос на 500~гПа не подтверждён, разд.~\ref{sec:limits}); к изученным областям и широтам; преимущественно к разрешённым ступеням лестницы масштабов, обе полосы которых не мельче 200~км (разд.~\ref{sec:construct}); к временным окнам 2017--2024~гг. Утверждение~\eqref{eq:inv} ничего не говорит о динамике: оно не касается переноса энергии между масштабами, направления каскада и прогностической полезности величины. Эти вопросы проверены отдельно, и результаты проверок приведены в разд.~\ref{sec:limits}. Далее термин <<региональный инвариант>> употребляется как сокращение утверждения~\eqref{eq:inv} с перечисленными границами --- в операциональном смысле традиции \cite{Puzachenko2010}.

Настоящая статья открывает серию из трёх работ, выполненных по единой программе и опирающихся на общий открытый репозиторий протоколов, кода и результатов. Первая работа вводит дескриптор межуровневой организации, проверяет его статус регионального инварианта и очерчивает границы его применимости. Вторая обращается к собственно географическому вопросу --- чем определяется географическая дифференциация инварианта: переходя от региональных областей к мозаичной карте, она атрибутирует его географию внешним ковариатам и разделяет вклад экзогенных граничных условий и состояния циркуляции. Третья предлагает интерпретирующую рамку, в которой инвариант понимается как характеристика установившейся организации циркуляции при закреплённом поле внешних условий, и формулирует её проверяемые следствия. Работы построены так, чтобы читаться независимо: ни одно утверждение настоящей статьи не опирается на результаты последующих, а содержательная интерпретация географии дескриптора отложена во вторую работу сознательно --- здесь фиксируется воспроизводимый рисунок.

Исследование построено по фальсификационной схеме: параметры процедуры, критерии и единственная статистика каждой проверки фиксировались до обращения к соответствующим данным, а ключевые утверждения подтверждались на выборках, не участвовавших в их получении. Полные протоколы и журнал отступлений открыты в репозитории (см. <<Доступность данных и кода>>).

\section{Конструирование дескриптора}
\label{sec:construct}

Чтобы проверить утверждение~\eqref{eq:inv}, нужна величина с тремя свойствами. Она должна измерять совпадение географии активности двух соседних уровней в пределах территории. Она должна быть нечувствительна к общей интенсивности циркуляции, потому что режимы существенно различаются общей интенсивностью циркуляции и любая мера, чувствительная к амплитуде, будет измерять эту разницу вместо организации. И она должна быть очищена от вклада самой процедуры разложения на уровни. Ниже описан материал, на котором такая величина строится, и три шага её построения.

\subsection{Материал}

Исходное поле --- относительный вихрь $\omega=\partial v/\partial x-\partial u/\partial y$ на поверхности 850~гПа, вычисленный по составляющим ветра реанализа ERA5 \cite{Hersbach2020,Soci2024} на сетке $0{,}25^\circ$ с шагом 6~ч в широтно-зависимой километровой метрике сферы. Вихрь подчёркивает вращательную структуру течения, из которой складываются циклоны, фронтальные волны и конвективные ячейки, и подавляет крупномасштабный адвективный фон. Уровень 850~гПа представителен для нижнетропосферной циркуляции, непосредственно взаимодействующей с подстилающей поверхностью.

Анализ ведётся в 39 областях размером $20^\circ$ широты $\times$ $40^\circ$ долготы. Двенадцать областей выбраны до начала вычислений так, чтобы охватить шесть контрастных циркуляционных режимов: тропические океаны, внетропические штормовые пути обоих полушарий, очаги глубокой конвекции над сушей, континентальные области умеренного и субтропического поясов (табл.~\ref{tab:regions}). Остальные 27 заданы формальным правилом, зафиксированным до получения данных: регулярная решётка ячеек $20^\circ\times40^\circ$ (пять широтных полос от $55^\circ$~с.\,ш. до $45^\circ$~ю.\,ш., долготные грани через $40^\circ$ от $20^\circ$~в.\,д.; 45 ячеек) с исключением ячеек, перекрывающихся с основными областями более чем на 20\,\% площади. Решёточные области используются только для проверки географии на составе выборки, не зависевшем от выбора автора (окна 2024~г.).

Для двенадцати основных областей привлечены три периода: 2017--2019~гг. (области R1--R4, 16 регион-окон, использованных при постановке метода), 2021--2022~гг. (области R5--R12, 32 регион-окна) и 2023--2024~гг. (все 12 областей, 48 регион-окон); каждая область в итоге располагает восемью сезонными окнами января--марта и июля--сентября. Периоды охватывают Ла-Нинью (2021--2022~гг.), Эль-Ниньо (2023--2024~гг.) и преимущественно нейтральные условия (2017--2019~гг.). Проверки кластеризации строятся на 48 основных регион-окнах; 32 окна 2021--2022~гг., охватывающие восемь из двенадцати областей, служат отложенной выборкой, не участвовавшей в построении дескриптора.

\begin{table}[tbp]
\centering\footnotesize
\caption{Двенадцать основных областей. Границы --- по протоколам получения данных; режим --- по классификации, принятой до начала вычислений. Периоды: I --- 2017--2019, II --- 2021--2022, III --- 2023--2024~гг.}
\label{tab:regions}
\begin{tabular}{@{}l>{\raggedright\arraybackslash}p{0.27\linewidth}>{\raggedright\arraybackslash}p{0.26\linewidth}>{\raggedright\arraybackslash}p{0.22\linewidth}l@{}}
\toprule
Код & Область & Границы & Циркуляционный режим & Периоды\\
\midrule
R1 & Тёплый бассейн запада тропической части Тихого океана & $10^\circ$~ю.\,ш.--$10^\circ$~с.\,ш.; $130$--$170^\circ$~в.\,д. & тропический океан & I, III\\
R2 & Северная Атлантика & $35$--$55^\circ$~с.\,ш.; $60$--$20^\circ$~з.\,д. & штормовой путь Северного полушария & I, III\\
R3 & Амазония & $15^\circ$~ю.\,ш.--$5^\circ$~с.\,ш.; $75$--$35^\circ$~з.\,д. & глубокая конвекция над сушей & I, III\\
R4 & Центральная Азия & $35$--$55^\circ$~с.\,ш.; $60$--$100^\circ$~в.\,д. & континентальная область умеренного пояса & I, III\\
R5 & Южно-Тихоокеанская зона конвергенции & $20^\circ$~ю.\,ш.--$0^\circ$; $180$--$140^\circ$~з.\,д. & тропический океан & II, III\\
R6 & Южная Атлантика & $55$--$35^\circ$~ю.\,ш.; $50$--$10^\circ$~з.\,д. & штормовой путь Южного полушария & II, III\\
R7 & Бассейн Конго & $15^\circ$~ю.\,ш.--$5^\circ$~с.\,ш.; $5$--$45^\circ$~в.\,д. & глубокая конвекция над сушей & II, III\\
R8 & Австралия & $35$--$15^\circ$~ю.\,ш.; $115$--$155^\circ$~в.\,д. & континентальная область субтропического пояса & II, III\\
R9 & Север Тихого океана & $35$--$55^\circ$~с.\,ш.; $180$--$140^\circ$~з.\,д. & штормовой путь Северного полушария & II, III\\
R10 & Индийский океан & $10^\circ$~ю.\,ш.--$10^\circ$~с.\,ш.; $60$--$100^\circ$~в.\,д. & тропический океан & II, III\\
R11 & Европа & $35$--$55^\circ$~с.\,ш.; $10^\circ$~з.\,д.--$30^\circ$~в.\,д. & континентальная область умеренного пояса & II, III\\
R12 & Юг Южной Америки & $40$--$20^\circ$~ю.\,ш.; $75$--$35^\circ$~з.\,д. & континентальная область субтропического пояса & II, III\\
\bottomrule
\end{tabular}
\end{table}

\subsection{Шаг 1. Локализация активности масштаба}

Волны разной длины знакопеременны и накладываются друг на друга. По самому полю вихря невозможно сказать, где в данный срок сосредоточена активность масштаба 200~км, а где --- масштаба 800~км. Поэтому поле разлагается по лестнице масштабов $\ell$ = 50, 100, 200, 400, 800 и 1600~км: для каждого $\ell$ строится изотропное гауссово сглаживание $\overline{\omega}_\ell$ с $\sigma=\ell/\sqrt{12}$ (дисперсионный эквивалент прямоугольного окна ширины $\ell$; ширина ядра в узлах сетки адаптируется к широтной метрике построчно), и разность соседних сглаживаний выделяет то, что живёт между двумя масштабами, --- октавный уровень
\begin{equation}
\label{eq:level}
D_i(x,y,t)=\overline{\omega}_{\ell_i}-\overline{\omega}_{\ell_{i+1}},\qquad i=1,\dots,5.
\end{equation}
Нулевой уровень --- остаток мельче первого масштаба, $D_0=\omega-\overline{\omega}_{50}$; всего получается шесть уровней $D_0,\dots,D_5$, от остатка мельче 50~км до полосы 800--1600~км. Такое разложение идейно соответствует выделению уровней организации по спектру у Кренке и соавторов \cite{Krenke2019} и переносит анализ в физическое пространство, с сохранением карты.

Уровень $D_i$ остаётся знакопеременным. Карту активности масштаба даёт модуль уровня, сглаженный на крупном краю его октавы (для нулевого уровня --- при $\sigma=\ell_1/\sqrt{12}$), --- огибающая
\begin{equation}
\label{eq:env}
E_i(x,y,t)=\overline{|D_i|}_{\ell_{i+1}}.
\end{equation}
Она показывает интенсивность вихревой активности данного масштаба в каждой точке, подобно тому как огибающая радиосигнала показывает громкость без слежения за каждым колебанием несущей. Огибающие как индикаторы локальной активности масштаба давно применяются в вейвлет-анализе турбулентности \cite{Farge1992}.

\subsection{Шаг 2. Совпадение рисунка карт}

После первого шага на каждый срок имеется шесть карт активности $E_0,\dots,E_5$, и вопрос гипотезы принимает конкретную форму: совпадает ли география активности соседних уровней? Размещается ли мезомасштабная активность там, где активен соседний крупный уровень? Искомая величина --- геометрическая ко-локализация двух карт. Карты предварительно логарифмируются: динамический диапазон вихря огромен, и логарифм выравнивает его, после чего сильные и слабые области входят в сравнение на равных. Для каждого 6-часового срока $t$ по узлам внутренней маски области (краевая полоса шириной $\min(2\ell_{\max};\ 0{,}25$ размера области$)$ исключена, чтобы в сравнение не попадали эффекты границ фильтрации) вычисляется пространственная ранговая корреляция логарифмов огибающих:
\begin{equation}
\label{eq:rho}
\rho_i(t)=\operatorname{corr}_{\mathrm{rank}}\bigl[\ln(E_{i-1}+\varepsilon),\ \ln(E_i+\varepsilon)\bigr],\qquad \varepsilon=10^{-12}\ \text{с}^{-1}.
\end{equation}
Корреляция берётся в ранговой форме как более устойчивая к выбросам; линейная корреляция логарифмов воспроизводит результаты с точностью $\rho=0{,}99$, так что выбор формы на выводы не влияет.

Для пяти пар соседних уровней получаем пять чисел на срок; медиана по ${\approx}360$ срокам сезонного окна (91~сут $\times$ 4) даёт профиль сопряжённости области $P=(\rho_1,\dots,\rho_5)$. Ступень 1 отвечает паре самых тонких полос (мельче 50; 50--100~км), ступень 2 --- паре (50--100; 100--200~км), ступень 3 --- (100--200; 200--400~км), ступень 4 --- (200--400; 400--800~км), ступень 5 --- паре самых крупных полос (400--800; 800--1600~км). Здесь нужно сразу договориться о доверии к ступеням. Оценки эффективного разрешения ERA5 расходятся: правило 4--7 шагов сетки \cite{Skamarock2004} даёт около 125--220~км, прямая спектральная оценка \cite{Bolgiani2022} --- 500--600~км в средних широтах и порядка 1000~км в тропиках. Полосы тоньше 200~км в любом случае в значительной мере определяются численной диссипацией и параметризациями модели. \emph{Разрешёнными} далее называются ступени 4--5, у которых обе полосы не мельче 200~км (умеренный порог по правилу 4--7 шагов). Все основные выводы дублируются контролем, ограниченным разрешёнными ступенями, а сопоставление с внешними источниками ведётся только по ним. Решающие проверки природы сигнала --- суррогатная поправка и воспроизведение географии в другом реанализе и в свободном расчёте --- от выбора порога не зависят: в них сравниваются поля, прошедшие один и тот же вычислительный тракт.

\subsection{Шаг 3. Вычитание вклада процедуры}

Соседние гауссовы полосы частично перекрываются, поэтому огибающие соседних уровней коррелируют уже в силу построения --- в том числе для случайного поля с тем же спектром, у которого нет негауссовой пространственной организации. Эту часть корреляции нужно измерить и вычесть. Для каждого регион-окна строится ансамбль из 199 фазовых суррогатов исходного поля ветра: на каждом сроке спектры $u$ и $v$ умножаются на общее случайное фазовое поле, так что амплитудный спектр каждой составляющей и среднее по области сохраняются точно, общая фаза сохраняет взаимную когерентность $u$ и $v$, а фазовая структура поля, в которой заключена его негауссова организация, разрушается. В силу линейности оператора вихря процедура тождественна фазовой рандомизации самого поля вихря с точным сохранением его пространственного спектра. Рандомизация выполняется на прямоугольнике области без оконной функции; краевые эффекты одинаковы для реального и суррогатного полей, проходящих один и тот же вычислительный тракт, и дополнительно подавляются внутренней маской. Вся цепочка операций повторяется на суррогатах, и основной формой дескриптора становится превышение над медианой суррогатного ансамбля, $P-\Psurr$ (далее --- профиль с суррогатной поправкой). Эта величина свободна от вклада перекрытия фильтров и измеряет негауссову организацию поля, которая в спектре отсутствует. Исходный профиль без поправки приводится параллельно во всех проверках.

Суррогатная поправка отвечает и на известную критику родственного коэффициента амплитудной модуляции: в \cite{CuiJacobi2025} показано, что тот отражает главным образом линейное взаимодействие со средним потоком. Здесь искомая часть сигнала выделяется контролем: поправка сравнивает наблюдаемую сопряжённость с полем той же спектральной структуры и удаляет весь вклад второго порядка по построению. Прямое сопоставление с самим коэффициентом амплитудной модуляции выполнено: на общем материале он статистически отличен от $P$ (ранговая корреляция карт $+0{,}14$) и не воспроизводит устойчивой региональной географии; предложенный дескриптор, таким образом, содержательно отличается от известной статистики.

Все параметры процедуры --- лестница масштабов, ширины фильтров, форма корреляции, размер краевой полосы, число суррогатов --- зафиксированы до начала вычислений и едины для всех областей, источников данных и проверок (табл.~\ref{tab:params}). Чувствительность к семейству фильтров, ширине сглаживания огибающей и отношению соседних масштабов после получения результатов не варьировалась; её систематическая проверка составляет отдельную методическую задачу.

\begin{table}[tbp]
\centering\small
\caption{Параметры процедуры оценивания (зафиксированы до начала вычислений).}
\label{tab:params}
\begin{tabular}{@{}p{0.34\linewidth}p{0.6\linewidth}@{}}
\toprule
Параметр & Значение\\
\midrule
Поле & относительный вихрь, 850~гПа, метрика сферы\\
Сетка и шаг по времени & $0{,}25^\circ$; 6~ч; окно ${\approx}360$ сроков (91~сут $\times$ 4)\\
Лестница масштабов $\ell$ & 50, 100, 200, 400, 800, 1600~км (октавы)\\
Фильтр уровня & изотропный гауссов, $\sigma=\ell/\sqrt{12}$, построчная адаптация к широтной метрике\\
Огибающая $E_i$ & $|D_i|$, сглаженный при $\sigma=\ell_{i+1}/\sqrt{12}$\\
Регуляризация логарифма & $\ln(E+\varepsilon)$, $\varepsilon=10^{-12}$~с$^{-1}$\\
Область корреляции & узлы внутренней маски области; краевая полоса $\min(2\ell_{\max};\ 0{,}25$ размера области$)$ исключена\\
Форма корреляции & ранговая (Спирмена); линейная форма даёт $\rho=0{,}99$ с ранговой\\
Единица вывода & регион-окно: медиана $\rho_i(t)$ по срокам окна\\
Суррогаты & фазовая рандомизация поля ветра, общее фазовое поле для $u$ и $v$; 199 реализаций на регион-окно; поправка по медиане ансамбля\\
\bottomrule
\end{tabular}
\end{table}

\section{Проверка гипотезы: устойчивость региональных профилей}
\label{sec:test}

Профиль сопряжённости монотонно убывает от мелких уровней к крупным во всех областях, а скорость убывания и уровень, на котором происходит основная потеря связи, различаются регионально (рис.~\ref{fig:profiles}). Тонкие линии на рисунке --- отдельные сезонные окна: зимние и летние, из разных лет и разных фаз ЭНСО. Их близость внутри каждой панели и различие между панелями --- то, что утверждение~\eqref{eq:inv} требует проверить количественно.

\begin{figure}[tbp]
\centering
\includegraphics[width=\linewidth]{figures/fig1_profiles.png}
\caption{Профили сопряжённости уровней $P$ для двенадцати основных областей (коды --- табл.~\ref{tab:regions}): тонкие линии --- отдельные сезонные окна, жирные --- средние по периодам 2017--2019, 2021--2022 и 2023--2024~гг. Ступени 1--5 отвечают парам соседних уровней: от (мельче~50; 50--100~км) до (400--800; 800--1600~км).}
\label{fig:profiles}
\end{figure}

Единица вывода --- регион-окно: медианное агрегирование по ${\approx}360$ срокам сводит каждое сезонное окно области к одному пятимерному вектору, так что внутриоконная временная зависимость сроков устраняется агрегированием и в статистику не входит. Дескрипторы стандартизуются покомпонентно (нулевое среднее и единичная дисперсия по выборке); статистика проверки --- разность медианных попарных евклидовых расстояний <<между регионами>> и <<внутри регионов>>, $\Delta=\operatorname{med}_{\mathrm{between}}-\operatorname{med}_{\mathrm{within}}$, в единицах стандартизованного пространства дескрипторов; она приводится всюду рядом с $p$ как размер эффекта. Нуль-модель --- 999 случайных переназначений региональных меток целым регион-окнам; $p=(1+\#\{\Delta^*\ge\Delta\})/1000$, поэтому минимально достижимое значение равно $p=0{,}001$. Перестановке подвергаются целые окна: межгодовая согласованность окон одного региона составляет само проверяемое свойство~\eqref{eq:inv}.

Региональная сигнатура подтверждается на 32 отложенных регион-окнах 2021--2022~гг., не участвовавших в построении дескриптора, с большим размером эффекта: $\Delta=+1{,}12$, $p=0{,}001$. Расстояния между профилями разных областей превышают расстояния между окнами одной области более чем на одну единицу стандартизованного пространства, и это сохраняется при смене сезона, независимого года и фазы ЭНСО (табл.~\ref{tab:summary}, строка~1). Ту же устойчивость показывает проверка на уровне целых регионов, не зависящая от вопроса об эффективном числе окон: корреляция географии профилей между двумя независимыми периодами ERA5 составляет 0,93 (по тем восьми областям, для которых доступны два периода).

Статистические ограничения схемы признаются явно. Перестановка региональных меток не моделирует пространственную корреляцию соседних областей, а окна одного региона разделяют режим и географию, поэтому число регион-окон (48 основных, 32 отложенных) --- верхняя оценка числа независимых единиц; консервативная нижняя --- число областей; усилением схемы были бы исключение целых регионов и блочный бутстреп по синоптическим отрезкам. Поправка на множественность проверок не вводилась: каждая строка табл.~\ref{tab:summary} отвечает отдельной, заранее заявленной гипотезе со своей, также заранее заявленной статистикой; значения на границе значимости ($p=0{,}026$--$0{,}033$) следует читать с обеими оговорками. На верхней ступени внутренняя часть области вмещает лишь несколько корреляционных радиусов, что ограничивает точность отдельных значений $\rho_5$; значимость при этом оценивается перестановками целых окон. Ключевые утверждения работы опираются на результаты с $p=0{,}001$ и большими размерами эффекта ($\Delta>1$), устойчивые к разумным поправкам эффективного объёма выборки.

\section{Методологические опасения и контроли}
\label{sec:controls}

Положительный результат разд.~\ref{sec:test} допускает четыре объяснения, в которых атмосфера участвует лишь косвенно. Каждое из них проверено отдельным контролем, заявленным до вычислений; сводка --- в табл.~\ref{tab:summary}.

\subsection{Опасение первое: разрешение реанализа}

Тонкие ступени лестницы лежат ниже эффективного разрешения реанализа, и их значения могли бы отражать свойства численной диссипации модели. Контроль, ограниченный разрешёнными ступенями 4--5, сохраняет региональную сигнатуру: $\Delta=+0{,}61$, $p=0{,}002$ (строка~7). Сопоставления с внешними источниками (разд.~\ref{sec:obs}) ведутся только по разрешённым ступеням.

\subsection{Опасение второе: спектр мощности}

Спектр мощности вихря сам по себе различает регионы (см. ниже), и региональные различия профиля могли бы быть пересказом спектра в других терминах. Проверка ведётся на регрессионных остатках по схеме с исключением по одному: из дескриптора удаляется часть, линейно предсказуемая по контрольному набору признаков, после чего тест кластеризации повторяется на остатках. Сигнатура сохраняется сверх полуфрактальных характеристик уровней --- прямых аналогов признаков выделения уровней в \cite{Krenke2019} ($\Delta=+0{,}65$, $p=0{,}001$), сверх скаттеринг-статистик второго порядка \cite{Cheng2024} ($\Delta=+0{,}30$, $p=0{,}026$) и сверх полного спектра мощности ($\Delta=+0{,}34$, $p=0{,}029$; строки 2--3). Две последние проверки находятся на границе значимости, и вопрос о спектре окончательно решается контролем равновеликими областями (разд.~\ref{sec:geom}): на них сверхспектральная кластеризация значима при $\Delta=+0{,}22$, $p=0{,}033$ для профиля с суррогатной поправкой и $\Delta=+0{,}35$, $p=0{,}007$ для исходного профиля (строка~6).

Методическая оговорка сохраняется: спектральные признаки сами по себе кластеризуют регионы сильнее профиля ($\Delta=+2{,}54$ на равновеликих областях). Утверждение работы состоит в наличии воспроизводимой информации сверх спектра, и именно оно подтверждено.

\subsection{Опасение третье: геометрия сетки}
\label{sec:geom}

Градусные области имеют широтно-зависимые километровые размеры, и часть межрегиональных различий могла быть картографическим отпечатком. Регрессионный контроль --- добавление модуля широты, шага сетки и ширины области в набор контрольных признаков --- сводил сверхспектральную кластеризацию к незначимой ($\Delta=+0{,}22$, $p=0{,}10$; строка~3) и оставлял вопрос открытым.

Контроль конструкцией решает его на уровне построения выборки. Те же двенадцать регионов заново нарезаны равновеликими областями $2000\times2800$~км (те же центры, индексные маски на исходной сетке), и вся последовательность операций, включая повторное построение фазовых суррогатов для каждого окна, повторена на 48 основных и 32 отложенных регион-окнах. Равновеликие области близки по площади к среднеширотным градусным и меньше экваториальных (5,6 против 9,9~млн~км$^2$); цель контроля --- равенство площадей между регионами. При такой нарезке региональная сигнатура сохраняется на основных и отложенных окнах ($\Delta=+1{,}26$ и $+1{,}56$, оба $p=0{,}001$; строка~5), география профилей при смене нарезки воспроизводится (объединённая ранговая корреляция <<градусы--километры>> 0,96; по отдельным регионам 0,70--1,00; рис.~\ref{fig:geometry}), а сверхспектральная кластеризация значима (строка~6). Расхождение двух контролей объяснимо: регрессионные ковариаты поглощали вместе с геометрией и физический межрегиональный сигнал, коррелированный с широтой; контроль построением выборки свободен от этого смешения.

\begin{figure}[tbp]
\centering
\includegraphics[width=\linewidth]{figures/fig3_geometry.png}
\caption{Контроль геометрии построением выборки: профили сопряжённости с суррогатной поправкой $P-\Psurr$ на градусных областях $20^\circ\times40^\circ$ (синий, <<degree>>) и на равновеликих областях $2000\times2800$~км (красный, <<km>>) для двенадцати регионов, 48 основных регион-окон 2023--2024~гг. Полосы --- размах по окнам; в заголовках панелей --- ранговая корреляция профилей двух нарезок по региону; по оси абсцисс --- полосы ступеней, км.}
\label{fig:geometry}
\end{figure}

Два дополнения снимают смежные вопросы о произвольности конструкции. Замена фиксированной октавной лестницы на лестницу, масштабированную локальным радиусом деформации, сохраняет региональную сигнатуру профиля с суррогатной поправкой ($\Delta=+1{,}60$, $p=0{,}001$): кластеризация не привязана к конкретному километражу ступеней. Вариация зональной протяжённости равновеликих областей на $\pm10\,\%$ меняет отношение сторон области при почти неизменной площади; профили каждого региона при этом воспроизводятся с $\rho\ge0{,}93$ (наименьшая по регионам ранговая корреляция 0,928), так что форма области не служит скрытым источником сигнала.

\subsection{Опасение четвёртое: наблюдательная сеть}
\label{sec:obs}

Реанализ восстанавливает состояние атмосферы моделью, усваивающей наблюдения, и плотность наблюдений неоднородна по земному шару. Возникает вопрос, не отражает ли региональная география сопряжённости размещение станций и спутниковых треков.

Первая проверка --- второй реанализ. MERRA-2 \cite{Gelaro2017} построен на другой модели общей циркуляции и другой системе усвоения; его сетка ($0{,}5^\circ\times0{,}625^\circ$) не разрешает полос мельче 200~км, поэтому лестница масштабов начинается с 200~км и сопоставление ведётся по разрешённым ступеням. Региональная сигнатура профиля воспроизводится в MERRA-2 с тем же размером эффекта, что и на отложенных окнах ERA5 ($\Delta=+1{,}15$, $p=0{,}001$; 48 регион-окон 2023--2024~гг.), а региональная география профиля по разрешённым ступеням совпадает с ERA5 с ранговой корреляцией 0,86 по двенадцати областям ($p=3\cdot10^{-4}$; строка~4). Оговорка о статусе этих чисел: протокол межреанализной проверки был заморожен для другой величины программы --- индекса невзаимности $A$ (разд.~\ref{sec:limits}), и именно для него получены $\Delta=+1{,}12$ и согласие географии 0,93; для профиля сопряжённости обе статистики строки~4 описательные, вычислены по тем же окнам после факта и в этом качестве приводятся. Проверка квазинезависима: оба реанализа усваивают в основном одни и те же наблюдения \cite{Hersbach2020,Gelaro2017}, поэтому она исключает артефакты конкретной системы усвоения и оставляет открытым вопрос о самой сети.

Решающая проверка --- расчёт без усвоения. В эксперименте \texttt{highresSST-present} проекта HighResMIP \cite{Haarsma2016} модель ECMWF-IFS-HR \cite{Roberts2018} интегрируется свободно при заданной наблюдаемой температуре поверхности океана и не получает атмосферных наблюдений. Использованы сезонные окна января--марта и июля--сентября 2014~г. --- последнего года эксперимента (протокол охватывает 1950--2014~гг.), поэтому сопоставление с окнами ERA5 2017--2024~гг. по построению климатологическое: сравниваются устойчивые географии. Если бы география сопряжённости была отпечатком сети, в таком расчёте она бы отсутствовала. По 39 областям география свободного расчёта воспроизводит географию ERA5 с ранговой корреляцией $\rho=0{,}71$ ($p=0{,}001$; рис.~\ref{fig:freerun}; строка~8). Оценивать эту величину следует относительно внутреннего предела метода --- корреляции между двумя независимыми периодами самого ERA5, равной 0,93 (восемь областей с двумя периодами): достигнуто около трёх четвертей возможного ($0{,}71/0{,}93=0{,}76$). Обе величины относятся к региональному уровню сопоставления; сопоставление глобальных мозаичных карт дескриптора составляет предмет отдельной работы и характеризуется собственными числами. Недостача ожидаема, поскольку свободный расчёт воспроизводит климатологию режимов, тогда как окна ERA5 относятся к конкретным годам. Сопоставление ведётся на общей сетке $0{,}5^\circ$ и только по разрешённым ступеням; при пессимистичной спектральной оценке разрешения ступень~4 частично попадает в серую зону, ступень~5 надёжна при обеих оценках.

\begin{figure}[tbp]
\centering
\includegraphics[width=\linewidth]{figures/fig4_free_running.png}
\caption{Сопоставление региональной географии по ERA5 (ось абсцисс) и по свободному расчёту ECMWF-IFS-HR без усвоения наблюдений (ось ординат) для 39 областей. Слева направо: профиль сопряжённости $P$; отвергнутый индекс невзаимности $A$ (показан для полноты: воспроизводимость индекса не оспаривается, отвергнуто его истолкование, разд.~\ref{sec:limits}); независимая информационная мера межуровневой связи; временная устойчивость карт активности. Пунктир --- линия равенства.}
\label{fig:freerun}
\end{figure}

Итог проверки: наблюдаемая региональная структура воспроизводится в свободной модельной интеграции и потому не может быть объяснена только устройством сети наблюдений или процедурой их усвоения. Полной независимости от всего, что модель разделяет с реанализом, один расчёт установить не может: общими остаются заданная температура океана, родственные параметризации и близкое эффективное разрешение. Желательное усиление --- ансамбль свободных реализаций и модель с иным динамическим ядром; расчёт CMCC-CM2-VHR4 доступен в том же архиве HighResMIP за тот же период.

\begin{table}[tbp]
\centering\small
\caption{Сводка проверок профиля сопряжённости. $\Delta$ --- размер эффекта теста кластеризации (разность медианных межрегиональных и внутрирегиональных расстояний в единицах стандартизованного пространства дескрипторов); $p$ --- перестановочный уровень значимости (999 переназначений региональных меток регион-окнам; минимум 0,001). Строки 1--4 --- 32 отложенных регион-окна 2021--2022~гг. и 48 основных 2023--2024~гг.; строки 5--6 --- контроль равновеликими областями. Строка~4 --- описательные статистики: протокол межреанализной проверки фиксировался для индекса~$A$ (разд.~\ref{sec:limits}).}
\label{tab:summary}
\begin{tabular}{@{}p{0.5\linewidth}p{0.45\linewidth}@{}}
\toprule
Проверка & Итог\\
\midrule
1. Региональная сигнатура: отложенные окна, независимые годы и фазы ЭНСО & $\Delta=+1{,}12$; $p=0{,}001$\\
2. Сверх полуфрактальных характеристик уровней \cite{Krenke2019}; сверх скаттеринг-статистик \cite{Cheng2024} & $\Delta=+0{,}65$; $p=0{,}001$ / $\Delta=+0{,}30$; $p=0{,}026$\\
3. Сверх полного спектра мощности & $\Delta=+0{,}34$; $p=0{,}029$; регрессионный контроль геометрии: $\Delta=+0{,}22$; $p=0{,}10$\\
4. Сигнатура профиля в MERRA-2; согласие географии ERA5--MERRA-2 по разрешённым ступеням (обе статистики описательные) & $\Delta=+1{,}15$; $p=0{,}001$; $\rho=0{,}86$; $p=3\cdot10^{-4}$\\
5. Контроль равновеликими областями: сигнатура; отложенные окна & $\Delta=+1{,}26$; $p=0{,}001$ / $\Delta=+1{,}56$; $p=0{,}001$\\
6. Контроль равновеликими областями: сверх спектра & профиль с поправкой $\Delta=+0{,}22$; $p=0{,}033$; исходный $\Delta=+0{,}35$; $p=0{,}007$\\
7. Только разрешённые ступени 4--5 & $\Delta=+0{,}61$; $p=0{,}002$\\
8. География в расчёте без усвоения наблюдений (39 областей) & $\rho=0{,}71$; $p=0{,}001$; предел метода 0,93\\
\bottomrule
\end{tabular}
\end{table}

\section{География: что различает инвариант}
\label{sec:geo}

\subsection{Карта}

Среднее значение профиля по разрешённым ступеням 4--5 меняется по 39 областям от 0,51 (ячейка над Индонезийским архипелагом) до 0,73 (тёплый бассейн запада тропической части Тихого океана; рис.~\ref{fig:map}); значения отдельных ступеней выходят за этот диапазон.

\begin{figure}[tbp]
\centering
\includegraphics[width=\linewidth]{figures/fig2_map.png}
\caption{География сопряжённости уровней: средние значения профиля по разрешённым ступеням 4--5 (пары полос 200--400/400--800 и 400--800/800--1600~км) для 39 областей, ERA5, 850~гПа; окна 2023--2024~гг. (решёточные области --- 2024~г.). Двенадцать областей выбраны содержательно, 27 --- формальным решёточным правилом.}
\label{fig:map}
\end{figure}

География читается содержательно, и её прочтение зависит от того, какая часть лестницы масштабов рассматривается. На разрешённых ступенях 4--5 (полосы 200--1600~км) высокие значения дают, с одной стороны, внетропические штормовые пути обоих полушарий (медианы по восьми сезонным окнам: Северная Атлантика 0,69, Южная Атлантика 0,67, север Тихого океана 0,65), с другой --- отдельные очаги глубокой конвекции: тёплый бассейн запада тропической части Тихого океана даёт наибольшее значение набора (0,73), бассейн Конго --- 0,69. Низкие значения принадлежат Амазонии (0,52), континентальным областям умеренного пояса (Европа 0,54, Центральная Азия 0,55) и Южно-Тихоокеанской зоне конвергенции (0,56). На тонких полосах (50--400~км) порядок иной и проще: три первых места занимают штормовые пути (север Тихого океана 0,37, Северная Атлантика 0,36, Южная Атлантика 0,35), два последних --- влажно-конвективные области (Индийский океан 0,26, Амазония 0,24); значения тонких полос приведены для профиля с суррогатной поправкой на равновеликих областях. Все региональные числа этого абзаца --- медианы по восьми сезонным окнам каждой области; на рис.~\ref{fig:map} показаны значения по окнам 2023--2024~гг., поэтому отдельные числа карты отличаются от приведённых во втором знаке.

Отсюда формулировка, принятая далее: тенденция <<штормовые пути --- выше, влажно-конвективные режимы --- ниже>> отчётливо выражена на тонких полосах и сохраняется как тенденция на разрешённых ступенях, где имеет явные исключения --- тёплый бассейн запада Тихого океана и бассейн Конго. Содержательное объяснение этого различия составляет предмет второй статьи серии; здесь фиксируется сам воспроизводимый рисунок.

\subsection{Два полюса: штормовые пути и конвекция}

Прочтение величины прямое. Высокое значение ступени означает, что мезомасштабная активность размещается там же, где активен соседний крупный уровень: по карте крупного уровня можно указать, где в этот срок сосредоточен мезомасштаб. Низкое значение означает, что та же пара уровней распределена почти безразлично и мезомасштабная организация в значительной мере отвязана от крупного поля. На тонких полосах два полюса этого различия заняты внетропическими штормовыми путями и влажно-конвективными областями; на разрешённых ступенях штормовые пути сохраняют высокие значения, тогда как влажно-конвективные области расходятся: Амазония и Южно-Тихоокеанская зона конвергенции дают низкие значения, тёплый бассейн запада Тихого океана и бассейн Конго --- высокие. Сравнение режимов здесь возможно потому, что величина нечувствительна к общей интенсивности циркуляции (разд.~\ref{sec:construct}): различие фиксируется в организации поля.

\subsection{Конго и Амазония}

Показательный случай --- два очага глубокой конвекции над сушей в близких широтах. Существенно для статуса примера: бассейны Конго и Амазонии --- единственные два региона глубокой конвекции над сушей в заранее заданном наборе двенадцати областей, так что пара не является результатом перебора вариантов; выбор для иллюстрации сделан после вычислений, но среди заранее заданных кандидатов выбора не было.

При сходном режиме увлажнения бассейны обнаруживают качественно различную архитектуру иерархии (рис.~\ref{fig:congo}). Различие сосредоточено на верхней ступени профиля: медиана по восьми сезонным окнам каждой области (Конго --- 2021--2024~гг., Амазония --- 2017--2019 и 2023--2024~гг.) составляет $\rho_5=0{,}70$ против $0{,}42$, причём размахи по окнам --- 0,57--0,81 и 0,36--0,53 соответственно --- не перекрываются. На четвёртой ступени различие медиан составляет 0,08, на пятой --- 0,28: контраст сосредоточен на верхней разрешённой ступени; значения ступеней ниже эффективного разрешения в аргументации не используются. Контраст устойчив к геометрии выборки: для профиля с суррогатной поправкой (четыре окна 2023--2024~гг.) контраст пятой ступени при переходе от градусных к равновеликим областям сужается с 0,242 до 0,136, сохраняя знак и 56\,\% величины (Амазония $0{,}120\to0{,}198$; Конго $0{,}363\to0{,}334$). Для контекста: у тропических океанических областей набора $\rho_5$ занимает промежуточные значения (0,53--0,68), так что контраст Конго--Амазония не является краем распределения: это различие двух сопоставимых по режиму территорий.

\begin{figure}[tbp]
\centering
\includegraphics[width=\linewidth]{figures/fig5_congo_amazon.png}
\caption{Бассейны Конго (вверху) и Амазонки (внизу), январь--март 2023~г. Слева и в центре --- карты активности крупного (800--1600~км) и мелкого (400--800~км) уровней на один срок; белые изолинии --- активность крупного уровня, одни и те же на обеих картах ряда. Справа --- совместное распределение пространственных рангов активности двух уровней по всем срокам окна; пунктир --- диагональ полного соответствия. В Конго мезомасштабная активность привязана к крупному полю, в Амазонии --- нет.}
\label{fig:congo}
\end{figure}

Содержательно это означает следующее. В бассейне Конго очаги мезомасштабной активности размещаются преимущественно там же, где активен соседний крупный уровень: по крупномасштабному полю можно указать, где окажется мезомасштаб. В Амазонии та же пара уровней распределена почти безразлично: мезомасштабная организация в значительной мере отвязана от крупного поля. Оговорка интерпретации: пара не свободна от неконтролируемых различий --- сезонность осадков, орографическое окружение и антропогенные изменения растительного покрова двух бассейнов различаются, --- поэтому пример демонстрирует различительную способность дескриптора; механизм различия здесь не устанавливается.

Для районирования это означает, что две территории, попадающие в один класс по стандартным характеристикам режима --- климатологии осадков, конвективной доступной потенциальной энергии и наклону спектра вихря, --- разделяются по устройству межуровневых связей, то есть по признаку, прямо относящемуся к задаче генерализации. В этом состоит непосредственная польза дескриптора для физико-географического районирования.

\section{Границы метода}
\label{sec:limits}

\subsection{Что измеряет величина}

Прежде чем перечислять, чего дескриптор не делает, нужно удостовериться, что он делает заявленное: измеряет, насколько размещение мелкого уровня предопределено соседним крупным. Для этого построена независимая мера --- невосстановимость размещения $u_i$. Для каждого срока ранги полей $\ln E_{i-1}$ и $\ln E_i$ по узлам внутренней маски разбиваются на $k=16$ равночастотных корзин; по двумерной гистограмме корзин вычисляется взаимная информация $I_i$, нормируемая на энтропию корзинного распределения ($H=\ln k$ при равночастотном разбиении), и
\begin{equation}
\label{eq:u}
u_i(t)=1-\frac{I_i(t)}{H};
\end{equation}
далее, как и для профиля, берётся медиана по срокам окна и среднее по разрешённым ступеням. Ранговое разбиение делает меру инвариантной к монотонным преобразованиям амплитуды, а обращение $1-I/H$ придаёт ей смысл невосстановимости: чем выше $u$, тем слабее размещение мелкого уровня предопределено соседним крупным. Связь профиля с этой мерой отрицательна и тесна: $\rho=-0{,}93$ (39 областей). При этом связь профиля с амплитудой полос слаба ($\rho=+0{,}28$ с логарифмом дисперсии разрешённых полос, 12 основных областей), а с временн\'ой устойчивостью карт активности --- со средней по узлам и разрешённым полосам автокорреляцией огибающей на лаге 6~ч --- практически отсутствует ($\rho=+0{,}01$, 39 областей). Обе величины --- статистики одного совместного распределения, так что проверка внутренняя; она устанавливает главное: профиль является сводкой именно того, насколько размещение мезомасштабной изменчивости предопределено соседним крупным уровнем, и не является тенью ни амплитуды, ни инерции поля.

\subsection{Чего она не измеряет}

Проверенные и отвергнутые истолкования сведены в табл.~\ref{tab:rejected}; они очерчивают область корректного использования величины.

Дескриптор не измеряет межмасштабный поток кинетической энергии: прямое сопоставление с потоком, вычисленным методом пространственного огрубления \cite{Germano1992,Aluie2018}, объясняет доли процента его изменчивости, а характеристики поля деформации информативнее на два порядка. Универсальной формы профиля, задаваемой средовыми ковариатами, нет: коллапса профилей к единой кривой не наблюдается. Вертикальной универсальности не обнаружено: на уровне 500~гПа региональная сигнатура статистически не подтверждается при вдвое меньшей выборке; корректная формулировка --- отсутствие подтверждения, поскольку различие не доказано.

Первоначально наряду с профилем предлагался второй дескриптор --- индекс невзаимности $A$, мера непереставимости статистических операторов агрегирования и детализации, которой приписывался смысл необратимости географического обобщения. Индекс воспроизводим (именно для него фиксировался протокол межреанализной проверки: MERRA-2 $\rho=0{,}93$, $p=1{,}2\cdot10^{-5}$; свободный расчёт: $\rho=0{,}90$), однако проверка содержания его не подтвердила: на синтетических полях с управляемой восстановимостью мелкого уровня по крупному индекс на неё не реагирует ($\rho=-0{,}32$ и $-0{,}35$ в двух прочтениях, тогда как независимая информационная мера отзывается корректно в обоих: 0,92--0,93 по модулю, со знаком, отвечающим ориентации меры в каждом прочтении) и монотонно следует за временн\'ой устойчивостью поля ($\rho=1{,}00$ на синтетике; $+0{,}53$ на реальных данных против $+0{,}04$ у информационной меры). Причина методическая: операторы оцениваются по скользящему окну, и чем медленнее меняется поле, тем систематичнее смещается результат. Индекс сохраняет статус воспроизводимой характеристики временн\'ой устойчивости поля и выведен из основного изложения.

Естественное ожидание --- что мера подчинённости мезомасштаба крупному полю говорит о качестве прогноза и об обоснованности статистической детализации --- проверено на трёх независимых постановках (табл.~\ref{tab:rejected}, последние строки), и ни одна не подтвердилась: в каждом случае более сильным предиктором оказалась простая характеристика фонового состояния. Скорость роста ошибки ансамблевого прогноза управляется бароклинностью ($\rho=+0{,}86$ с модулем широты), мезомасштабная ошибка анализа --- амплитудой поля, а география потери мезомасштаба в нейросетевых моделях погоды \cite{Lam2023,Bi2023,Price2025,Rasp2024} предсказывается спектральной характеристикой лучше, чем дескриптором. Отрицательный итог следует читать буквально: на данных, собранных неоднородной наблюдательной сетью, и на моделях, целевые функции которых уничтожают мезомасштабную структуру, отсутствие связи не позволяет заключить ни о наличии, ни об отсутствии прикладного содержания величины. Для продвижения нужны наблюдения, однородные по плотности, и модели, сохраняющие мезомасштабную изменчивость; протоколы этих проверок и сопутствующие результаты приведены в репозитории.

\begin{table}[tbp]
\centering\small
\caption{Проверенные и отвергнутые интерпретации и применения. Первые шесть строк --- проверки содержания величины; последние три --- прикладные гипотезы.}
\label{tab:rejected}
\begin{tabular}{@{}p{0.5\linewidth}p{0.45\linewidth}@{}}
\toprule
Проверяемое положение & Итог проверки\\
\midrule
Индекс невзаимности $A$ измеряет необратимость обобщения & отвергнуто (следует за временн\'ой устойчивостью поля)\\
Дескрипторы отслеживают межмасштабный поток кинетической энергии \cite{Germano1992,Aluie2018} & отвергнуто (характеристики поля деформации информативнее на два порядка)\\
Форма профилей подчиняется универсальному закону от средовых ковариат & отвергнуто (коллапса к единой кривой нет)\\
Дескрипторы вертикально универсальны (перенос на 500~гПа) & не подтверждено (уровнеспецифичность)\\
Дескрипторы информативны для замыкания влагобюджета реанализа \cite{Trenberth1991,Mayer2021} & отвергнуто\\
Режимная связь с осадками формализуема как уравнение <<источник--сток>> \cite{Masunaga2012,PetersNeelin2006} & не установлено (характерное время ${\sim}2$~сут частично воспроизводится инерцией процедуры оценивания)\\
Дескрипторы предсказывают скорость роста ошибки ансамблевого прогноза & не подтверждено на имеющихся данных; сильнее предиктор --- бароклинность ($\rho=+0{,}86$ с модулем широты)\\
Дескрипторы предсказывают уровень мезомасштабной ошибки анализа & не подтверждено на имеющихся данных; связь исчезает после контроля амплитуды\\
Дескрипторы предсказывают потерю мезомасштаба в моделях машинного обучения \cite{Lam2023,Bi2023,Price2025,Rasp2024} & не подтверждено на имеющихся данных; сильнее предиктор --- спектральная характеристика\\
\bottomrule
\end{tabular}
\end{table}

Утверждение об инвариантности, таким образом, ограничено полем, уровнем, масштабами и периодом, перечисленными в разд.~\ref{sec:hyp}; за их пределами --- 500~гПа, другие переменные, прогностические применения --- свойства величины либо не подтверждены, либо отвергнуты. Работа не утверждает, что найден универсальный инвариант атмосферной динамики, самостоятельная динамическая переменная или мера межмасштабного переноса.

\section{Заключение}
\label{sec:concl}

Предложен инструмент измерения сопряжённости иерархических уровней атмосферной циркуляции --- профиль сопряжённости октавных уровней $P$ по полю относительного вихря на 850~гПа. Его построение полностью описано операционально, все параметры зафиксированы до вычислений, а ранговая форма корреляции представляет собой техническую деталь ($\rho=0{,}99$ с линейной формой); новизна работы состоит в свойствах измеряемого объекта.

Профиль удовлетворяет заранее сформулированному утверждению условной инвариантности~\eqref{eq:inv}: региональная сигнатура сохраняется при смене сезона, независимого года, фазы ЭНСО ($\Delta=+1{,}12$, $p=0{,}001$ на отложенных выборках), источника данных (MERRA-2: $\Delta=+1{,}15$, $p=0{,}001$; география $\rho=0{,}86$) и схемы разбиения территории, и не сводится к спектру мощности (сверхспектральная кластеризация на равновеликих областях: $\Delta=+0{,}22$, $p=0{,}033$ для профиля с суррогатной поправкой). Региональная география воспроизводится в свободном модельном расчёте, не усваивающем атмосферных наблюдений ($\rho=0{,}71$ при внутреннем пределе метода 0,93; региональный уровень сопоставления, 39 областей), поэтому наблюдаемая структура не может быть объяснена только устройством наблюдательной сети или процедурой усвоения.

География дескриптора систематична и масштабно-зависима: на тонких полосах первые места занимают внетропические штормовые пути, последние --- влажно-конвективные области; на разрешённых ступенях та же тенденция сохраняется с исключениями --- наибольшее значение набора принадлежит тёплому бассейну запада тропической части Тихого океана, высокое значение --- бассейну Конго. Бассейны Конго и Амазонии, единственные два очага глубокой конвекции над сушей в заранее заданном наборе регионов, различаются архитектурой иерархии с неперекрывающимися размахами по окнам ($\rho_5$: 0,57--0,81 против 0,36--0,53), и контраст сохраняется при контроле равновеликими областями. Дескриптор различает территории, близкие по стандартным характеристикам режима, и в этом качестве продолжает линию инвариантов организации геосистем --- от выделения уровней к измерению их взаимодействия.

Границы применимости установлены прямыми проверками: дескриптор не измеряет межмасштабный поток кинетической энергии; вертикальная универсальность не подтверждена; три прикладные гипотезы не подтверждены на имеющихся данных, причём в каждом случае более сильным предиктором оказалась простая характеристика фонового состояния; первоначально предлагавшийся индекс невзаимности сохраняет статус воспроизводимой характеристики временн\'ой устойчивости поля, а приписанное ему истолкование отвергнуто.

Два направления дальнейшей работы следуют из результатов непосредственно. Первое --- физико-географическое районирование по устройству межуровневых связей: величина даёт координату, дополняющую климатологию режима и различающую территории там, где стандартные характеристики их не разделяют. Второе --- диагностика воспроизведения мезомасштабной организации в климатических моделях: поскольку география величины воспроизводится свободным расчётом, она задаёт наблюдаемую цель, не зависящую ни от усвоения наблюдений, ни от абсолютной амплитуды поля; такая оценка в настоящей работе не выполняется и остаётся предложением.

\bigskip
\noindent\textbf{Финансирование.} Исследование выполнено без внешнего финансирования.

\smallskip
\noindent\textbf{Конфликт интересов.} Автор заявляет об отсутствии конфликта интересов.

\smallskip
\noindent\textbf{Доступность данных и кода.} Исходные данные --- открыто доступные продукты сторонних поставщиков и в репозитории автора не дублируются: поля ERA5 доступны через Climate Data Store (Copernicus Climate Change Service), данные MERRA-2 --- через NASA GES DISC, модельные расчёты HighResMIP --- через Earth System Grid Federation; доступ к первым двум требует бесплатной учётной записи. Код анализа, скрипты получения исходных данных, полные протоколы проверок с заранее фиксированными критериями, журнал отступлений и производные табличные результаты, достаточные для воспроизведения всех рисунков и статистических проверок работы, открыты в репозитории \url{https://github.com/Theclimateguy/Shroedinger} (версия v6.2; постоянный архив версий --- Zenodo, DOI 10.5281/zenodo.19565769). Маршрут воспроизведения статьи --- от получения данных до каждой строки табл.~\ref{tab:summary} --- описан в файле \texttt{reproducibility/article1\_README.md} и автоматизирован скриптом \texttt{reproducibility/reproduce\_article1.py}; там же приведены контрольные суммы ожидаемых артефактов и проверочный маршрут на одном регион-окне, не требующий массовой загрузки данных. Описательные величины, не входившие в зафиксированные тесты (региональные порядки на разрешённых и тонких ступенях, контраст Конго--Амазония, корреляции проверки содержания, география профиля ERA5--MERRA-2 из строки~4), пересчитываются скриптом \texttt{clean\_experiments/verify\_article1\_descriptives.py}.

\printbibliography[title={Список литературы}]

\end{document}
